\section{A framed category and its Artin action}
\label{sec:category}

We construct a finite graded algebra and a homotopy action of the Artin
group on complexes of its projectives. The construction uses a fusion
category to realize the nonintegral entries of the geometric Coxeter
form.  This is a framed variant of the fusion-category zigzag construction
and Artin action of Heng--Licata~\cite[Sections~4--5]{HengLicata2024}.
We use a full Temperley--Lieb factor for each edge with finite label
at least four, its string
object as edge label, and an additional frame.  These choices differ
from their default labels; compare their Remarks~4.4 and~5.4.
We give the algebra and its bimodule homotopies for these conventions.

\subsection{The geometric form and a frame}

Adjoin a vertex $o$, called the frame, to the original finite set $S$,
and put $\widehat S=S\sqcup\{o\}$. For distinct original vertices set
\[
 b_{ij}=\begin{cases}
 2\cos(\pi/m_{ij}),&m_{ij}<\infty,\\
 2,&m_{ij}=\infty,
 \end{cases}
 \qquad b_{io}=b_{oi}=1\quad(i\in S).
\]
Let $C$ be the symmetric matrix indexed by $\widehat S$ with diagonal
entries $2$ and off-diagonal entries $-b_{ij}$. An edge means a pair with
$b_{ij}>0$; in particular a label $2$ gives no edge. The frame is not an
Artin generator. If $T\subseteq\widehat S$, write $C_T$ for the
corresponding principal submatrix, and write $\partial T$ for the vertices
outside $T$ adjacent to a vertex of $T$. This boundary includes the frame
when appropriate.

We use the standard geometric representation of Coxeter systems: on
root-coefficient vectors the simple reflection at $i$ is
\begin{equation}
 s_i p=p-(Cp)_i\varepsilon_i,
 \label{eq:category-reflection}
\end{equation}
where $\varepsilon_i$ is the $i$th coordinate vector. Its restriction to
the coordinates of a parabolic subsystem is its geometric
representation. The standard Coxeter theorems identify spherical
$T\subseteq S$ with the positive definite matrices $C_T$
\cite[Chapter~V, \S4, no.~8, Theorem~2]{Bourbaki2002}. We also use
the root-sign criterion for descents and the fact that reduced
expressions are related by braid moves; these are Coxeter-system
theorems, without a finiteness assumption on the ambient system
\cite[Chapter~IV, \S1; Chapter~V, \S4, nos.~4 and~6]{Bourbaki2002}.
Whenever the chamber covering or a
longest element is used below, the subsystem concerned has first been
shown to be finite.

\begin{lemma}
\label{lem:cartan-positive-null}
For a spherical subset $T\subseteq S$, the inverse $C_T^{-1}$ is
entrywise nonnegative, and is strictly positive within each connected
diagonal block. Prescribing the coordinates of $p$ outside $T$ therefore
determines a unique vector harmonic on $T$, by which we mean
$(Cp)|_T=0$. Nonnegative prescribed coordinates give a nonnegative
harmonic extension.

More generally, let $G$ be a nonempty connected vertex set. If
$v_i>0$ for $i\in G$ and $C_Gv\geq0$ coordinatewise, then either
$C_G$ is positive definite or $C_Gv=0$.
\end{lemma}

\begin{proof}
Write $C_T=2I-B_T$, where $B_T$ is symmetric and entrywise nonnegative.
Positive definiteness implies that the largest eigenvalue of $B_T$ is
less than $2$. The Perron--Frobenius bound for a nonnegative symmetric
matrix then gives spectral radius less than $2$, so
\[
 C_T^{-1}=\frac12\sum_{n\geq0}(B_T/2)^n.
\]
This is nonnegative, and a path between any two vertices in the same
block gives a strictly positive entry in some power. Harmonic extension
is explicitly
\[
 p_T=C_T^{-1}B_{T,\widehat S\setminus T}
                 p_{\widehat S\setminus T},
\]
which proves the first assertions, including uniqueness.

For the last assertion, every real vector $x$ on $G$ satisfies
\begin{equation}
 x^{\mathsf t}C_Gx
 =\sum_{\substack{i<j\\i,j\in G}}b_{ij}v_iv_j
       \left(\frac{x_i}{v_i}-\frac{x_j}{v_j}\right)^2
   +\sum_{i\in G}\frac{(C_Gv)_i}{v_i}x_i^2.
 \label{eq:category-scaled-form}
\end{equation}
If at least one coordinate of $C_Gv$ is positive, vanishing of this
quadratic form forces all $x_i/v_i$ to be equal by connectedness, and
then forces their common value to be zero. The form is therefore
positive definite. If none is positive, the assumed inequalities say
exactly that $C_Gv=0$.
\end{proof}

\subsection{Fusion labels}

The ingredients are the Jones--Wenzl projectors and the spherical
quotient by negligible maps.  We recall their construction to fix the
normalization; see Wenzl~\cite{Wenzl1987},
Barrett--Westbury~\cite[Theorem~2.9]{BarrettWestbury1999}, and the
Temperley--Lieb account in Chen~\cite[Sections~5.1--5.4]{Chen2014}.
The finite semisimplicity below uses the specific projector decomposition,
in addition to the general spherical quotient construction.

\begin{proposition}
\label{prop:fusion-construction}
There is a finite semisimple spherical $\C$-linear tensor category
$\mathscr C$ with strictly positive real dimensions, and, for every
ordered edge $(s,t)$ of $\widehat S$, an object $D^1_{st}$ of dimension
$b_{st}$, with $D^1_{ts}$ its dual, having the following properties.
For every finite edge and simple objects $L,L'$,
\[
 \dim_\C\Hom_{\mathscr C}(L,D^1_{st}\otimes L')\leq1.
\]
On a label-$3$ edge or a frame edge one may take $D^1_{st}=\mathbf1$;
on an infinite edge one may take $D^1_{st}=\mathbf1\oplus\mathbf1$.
\end{proposition}

\begin{proof}
For each individual finite edge of label $m\geq4$, use an independent
Temperley--Lieb factor with circle value
\[
 d_m=2\cos(\pi/m),\qquad
 [j]=\frac{\sin(j\pi/m)}{\sin(\pi/m)}.
\]
Here is its construction. Begin with the pivotal diagram category whose
objects are nonnegative tensor powers of a self-dual string and whose
morphisms are $\C$-linear combinations of planar pairings. Composition
is stacking, a closed circle is replaced by $d_m$, and tensor product
is juxtaposition. Closing a diagram gives its spherical trace. Take the
additive idempotent completion, quotient by the ideal of negligible
maps, and complete additively and idempotently again. A map $f:X\to Y$
is negligible if $\operatorname{tr}(gf)=0$ for every $g:Y\to X$.
Cyclicity proves stability under composition, and partial trace proves
stability under tensoring; hence this is a tensor ideal.

For completeness, the idempotents and decompositions in this quotient
can be obtained directly. Let $f_0$ and $f_1$ be the respective identity
diagrams, and define, for $1\leq j\leq m-2$,
\begin{equation}
 f_{j+1}=p-\frac{[j]}{[j+1]}p e_jp,
 \qquad p=f_j\otimes\id,
 \label{eq:category-jw-recursion}
\end{equation}
where $e_j$ is the unnormalized cap--cup at the last two positions.
All denominators in this recursion are nonzero. Inductively $f_j$ is
idempotent, kills every elementary turnback on either side, and its
partial trace over the last string is
\begin{equation}
 \operatorname{ptr}(f_j)=\frac{[j+1]}{[j]}f_{j-1}
 \qquad(1\leq j\leq m-1).
 \label{eq:category-jw-partial-trace}
\end{equation}
Indeed the partial trace of the two terms defining $f_j$ is
$d_m f_{j-1}-([j-1]/[j])f_{j-1}^2$, which is the right side of
\eqref{eq:category-jw-partial-trace}; the case $j=1$ starts the
induction. Sandwiched between the last cap--cups, $p$ consequently gives
$([j+1]/[j])(f_{j-1}\otimes e)$, where $e$ is the last cap--cup.
Since $p$ already factors through $f_{j-1}$ on the first $j-1$ strings,
this identity shows that the subtracted term in
\eqref{eq:category-jw-recursion} is a projection, proves idempotence
of $f_{j+1}$, and proves its turnback cancellation. The last cap and cup
factor that projection through $f_{j-1}$; their opposite composite is
the nonzero scalar $[j+1]/[j]$. Thus its image is isomorphic to
$f_{j-1}$, giving
\begin{equation}
 f_j\otimes f_1\cong f_{j+1}\oplus f_{j-1}
 \qquad(1\leq j<m-1).
 \label{eq:category-string-fusion}
\end{equation}
Closing \eqref{eq:category-jw-partial-trace} gives
$\operatorname{tr}(f_j)=[j+1]$.

A planar pairing between the corners of $f_i$ and $f_j$ has a turnback
unless $i=j$ and it is the straight pairing. Consequently these corners
vanish for $i\ne j$ and are one-dimensional for $i=j$. The object
$f_{m-1}$, whose trace is $[m]=0$, is therefore negligible. The objects
$f_0,\ldots,f_{m-2}$ survive, are simple, and have positive dimensions
$[1],\ldots,[m-1]$. The decompositions
\eqref{eq:category-string-fusion}, with $f_{m-1}$ now zero, generate
every tensor power from these objects. Thus the quotient is finite
semisimple. Reflecting the diagrams gives the same fusion rule on the
other side: the turnback-killing idempotent with straight coefficient
one is unique, as follows by multiplying any two such idempotents.
Duality and the spherical trace descend to the quotient and its
completions. This proves all the stated properties of a single factor.

Take the tensor product of these finitely many categories, with simple
objects the tuples of simple objects in the factors and tensor product
and Hom spaces given factorwise. If there are no such factors, take
finite-dimensional vector spaces. For an edge of label $m\geq4$, use
the string object $f_1$ in its factor and units in the other factors.
Tensoring by it has the path-$A_{m-1}$ multiplicities in that factor and
fixes the others, proving the multiplicity bound. The choices of units
and two units cover the remaining edges. All dimensions are products
of positive dimensions. Choose dual objects for opposite orientations
and use the opposite pivotal evaluations on the two orientations.
\end{proof}

Fix representatives $\operatorname{Irr}(\mathscr C)$ of the simple
objects, and put
\[
 w_L=\dim_{\mathscr C}L>0
 \qquad(L\in\operatorname{Irr}(\mathscr C)).
\]
The tensor category is not required to be symmetric. In what follows
the labels are always tensored on the right, and evaluation of an edge
and its dual occurs on the left of those labels.

\subsection{The graded algebra and its pairing}

Ordinary zigzag algebras and their graded-symmetric pairing are developed
by Huerfano--Khovanov~\cite[Section~3, Proposition~1]{HuerfanoKhovanov2001}.
The following projective category is the basic-algebra realization of
the fusion-labelled version; compare
\cite[Section~4.4, Proposition~4.13]{HengLicata2024}.

For $s\in\widehat S$ and $X\in\mathscr C$ introduce an additive
projective symbol $P_{s,X}$. Besides $D^1_{st}$ on edges, set
$D^0_{ss}=D^2_{ss}=\mathbf1$ and set every unspecified $D^a_{st}$ to
zero. Define the length-$a$ morphisms by
\begin{equation}
 \Hom^a(P_{s,X},P_{t,Y})
   =\Hom_{\mathscr C}(X,D^a_{st}\otimes Y).
 \label{eq:category-projective-hom}
\end{equation}
The length-zero units act as identities. For two length-one maps,
composition is obtained by first tensoring their label maps and then
applying
\[
 D^1_{st}\otimes D^1_{ts}\longrightarrow\mathbf1=D^2_{ss}
\]
on a backtrack. All other products of two positive lengths are zero.
For opposite directions use the left and right pivotal evaluations.
This is associative: the assertion involving a length-zero map is
functoriality, and both products of any three positive lengths are
zero.

An unfolded vertex is a pair $\alpha=(s,L)$ with $L$ simple. For a set
of colors $T$ write $\widetilde T=T\times\operatorname{Irr}(\mathscr C)$.
Let $Z$ be the finite-dimensional basic graded algebra with
\begin{equation}
 e_\alpha Z e_\beta=\Hom(P_\alpha,P_\beta),
 \label{eq:category-path-order}
\end{equation}
where multiplication is in path order. With this convention
$P_\alpha=Ze_\alpha$ is a left projective and
$\Hom_Z(Ze_\alpha,Ze_\beta)=e_\alpha Ze_\beta$ as required. Put
$Q_\alpha=e_\alpha Z$ for the right projective.
Let $U_\alpha$ be the simple top of $Q_\alpha$, concentrated in
cochain and internal degrees zero. Decomposing $X$ into
simples expresses $P_{s,X}$ as a sum of the $P_\alpha$'s. The algebra
has only lengths $0,1,2$; length zero is the semisimple diagonal, and
length two is one-dimensional on each unfolded diagonal. Write
$t_\alpha$ for the length-two token corresponding to $\id_L$ when
$\alpha=(s,L)$.

\begin{lemma}
\label{lem:category-frobenius}
The functional $\tau:Z\to\C$ defined by
\[
 \tau(t_{s,L})=w_L,
 \qquad \tau(Z^0\oplus Z^1)=0
\]
is a symmetric Frobenius trace of degree $-2$. Thus
$\theta_\alpha=t_\alpha/w_L$ has trace one. The radical of $Z$ is
$Z^{>0}$, and the socle of each indecomposable left or right projective
is its token subspace.

For every finite graded bimodule $M$, the adjoint dual has a natural
bimodule identification
\begin{equation}
 M^\vee:=\Hom_{Z\text{-left}}(M,Z)
       \cong D_\C(M)\langle2\rangle,
 \label{eq:category-adjoint-duality}
\end{equation}
where $D_\C$ is graded vector-space duality and $\langle a\rangle$
raises degrees by $a$. It is an involutive contravariant equivalence,
fixes the regular bimodule, and satisfies
\begin{equation}
 (P_\alpha\otimes_\C Q_\beta)^\vee
       \cong(P_\beta\otimes_\C Q_\alpha)\langle-2\rangle.
 \label{eq:category-projective-dual}
\end{equation}
In particular this duality preserves additive radicals.
\end{lemma}

\begin{proof}
The length-zero and length-two spaces pair perfectly. For length one,
duality identifies maps $L'\to D^1_{ts}\otimes L$ with maps
$D^1_{st}\otimes L'\to L$. Pairing these against
$L\to D^1_{st}\otimes L'$ by the semisimple trace is nondegenerate:
on every simple summand the ordinary matrix trace is multiplied by its
nonzero categorical dimension. This is exactly the backtrack pairing.
Its symmetry is the pivotal partial-trace identity. Namely, the trace
on $L$ of the first backtrack equals the trace on
$D^1_{st}\otimes L'$ of the opposite composite; taking partial trace
in $D^1_{st}$ gives the second pivotal cap and then the trace on $L'$.
Hence $\tau(ab)=\tau(ba)$. This is ordinary symmetry; it has no sign
from internal degree.

The positive-length ideal is nilpotent and the quotient is the
semisimple diagonal, so it is the radical. Tokens are annihilated by
this radical. A nonzero length-one element has a nonzero product with
a suitable opposite arrow by the perfect pairing, and a nonzero
length-zero element has a nonzero product with its token. Therefore
the token is the entire socle of each projective.

For \eqref{eq:category-adjoint-duality}, send $f$ to $\tau\circ f$.
A homogeneous $f$ of degree $k$ gives a functional of degree $k-2$,
which explains the shift $\langle2\rangle$. Nondegeneracy of the trace
gives the usual Frobenius identification for every finite left module.
It respects both bimodule actions: with
$(a f b)(m)=f(ma)b$ one has
\[
 \tau(f(ma)b)=\tau(bf(ma))=\tau(f(bma)),
\]
which is the action $(a\lambda b)(m)=\lambda(bma)$ on the vector dual.
The shifted vector dual squares naturally to the identity, so this is
a contravariant equivalence. Such an equivalence preserves the
additive radical, characterized by invertibility of $1-gf$ for all
opposite maps $g$.

The trace identifies
$D_\C(P_\alpha)=Q_\alpha\langle-2\rangle$ and
$D_\C(Q_\beta)=P_\beta\langle-2\rangle$. Vector duality reverses
the factors of an external tensor product; combining its shift $-4$
with the shift $+2$ in \eqref{eq:category-adjoint-duality} proves
\eqref{eq:category-projective-dual}. The same trace identifies
$D_\C(Z)\langle2\rangle$ with $Z$.
\end{proof}

\subsection{Finite Coxeter unfoldings}

Fusion unfolding replaces labelled vertices by their simple constituents;
see Heng--Licata~\cite[Section~8]{HengLicata2024}.
For the chosen factors we prove the positivity and descent properties
that will be needed for finite-type braid and termination arguments.

Define a symmetric matrix $\widetilde C$, indexed by unfolded vertices,
with diagonal $2$ and off-diagonal entries
\[
 \widetilde C_{\alpha\beta}
       =-\dim_\C(e_\alpha Z^1e_\beta).
\]
Symmetry follows also from the perfect edge pairing. Define
\begin{equation}
 (\phi p)_{s,L}=w_Lp_s.
 \label{eq:category-phi}
\end{equation}
Fusion and duality imply
\begin{equation}
 \widetilde C\phi p=\phi Cp.
 \label{eq:category-intertwining}
\end{equation}
Indeed, for fixed $L$ and an edge object $D$, the weighted sum of
$\dim\Hom_{\mathscr C}(L,D\otimes L')$ over $L'$ is
$\dim(D)w_L$: adjunction identifies these multiplicities with those
of $L'$ in $D^*\otimes L$, whose dimension is $\dim(D)w_L$.
Equation~\eqref{eq:category-intertwining} also holds on any color
submatrix.

We write $D_L(w)=\{s:\ell(sw)<\ell(w)\}$ and
$D_R(w)=\{s:\ell(ws)<\ell(w)\}$ for the left and right descent sets,
where $s$ ranges over the simple generators of the Coxeter system
in question.

\begin{lemma}
\label{lem:finite-unfolding}
If $J\subseteq S$ is spherical, then $\widetilde C_{\widetilde J}$
is positive definite and simply laced, and defines a finite Coxeter
system. Replacing each original letter $s_i$, $i\in J$, by the product
of the commuting reflections in all vertices of its color gives a
lift of words to this system. For every such word representing $w$,
its lifted product $\widetilde w$ has
\begin{equation}
 D_L(\widetilde w)=\widetilde{D_L(w)},\qquad
 D_R(\widetilde w)=\widetilde{D_R(w)}.
 \label{eq:category-descents}
\end{equation}
Reduced expressions lift to reduced expressions, and a reduced
expression for the longest element of $W_J$ lifts to a reduced
expression for the longest element of the unfolded system.
\end{lemma}

\begin{proof}
The positive vector
$v=\phi(C_J^{-1}\mathbf1)$ satisfies
$\widetilde C_{\widetilde J}v=\phi\mathbf1>0$.
Applying the scaled quadratic-form identity
\eqref{eq:category-scaled-form} to each connected component proves
positive definiteness of the entire unfolded matrix. In particular
this proves positivity on directions outside the image of $\phi$.
An infinite edge cannot lie in $J$, since its two-by-two Cartan matrix
is singular. The multiplicity assertion of
Proposition~\ref{prop:fusion-construction} therefore makes every
off-diagonal unfolded entry either $0$ or $-1$. The finite Coxeter
criterion applies.

There are no edges between different vertices of a single color.
Their reflections commute, and
\eqref{eq:category-intertwining} shows that their product on a folded
vector is the original reflection
\eqref{eq:category-reflection}. Choose $p$ with $C_Jp>0$. Then
$\phi p$ lies in the interior of a full unfolded fundamental chamber,
since every row of $\widetilde C_{\widetilde J}\phi p$ is positive.
For any word, intertwining identifies its image with $\phi(wp)$.
The row at an unfolded vertex $\alpha$ on this image tests the sign
of $\widetilde w^{-1}\alpha$, while the corresponding original row
tests the sign of $w^{-1}\alpha_i$.

To see why this is a complete sign test, a root of either finite
system has either all nonnegative or all nonpositive coefficients in
the simple-root basis. Pairing it with one chamber-interior point is
therefore strictly positive or strictly negative respectively. No
assumption about the dimension of the folded subspace or about its
orbit meeting every unfolded chamber is needed. The intertwining
identity multiplies the original row by $w_L>0$, proving the equality
of left descent sets. Applying the same argument to reversed words
proves the equality of right descent sets. The lifted product depends
only on the original group element as well: a word for the identity
fixes $\phi p$, and the stabilizer of a chamber-interior point in the
finite unfolded reflection group is trivial
\cite[Chapter~V, \S3, no.~3]{Bourbaki2002}.

Suppose an original reduced expression is extended by an ascent
letter. Every vertex of its color is an ascent of the lifted prefix
by \eqref{eq:category-descents}. After one reflection of that color,
the others remain ascents, because they are orthogonal to it.
Consequently the whole block is a sequence of successive ascents.
Induction proves reducedness of the lift. Finally the longest
original element has every original simple reflection as a descent.
Its lift has every unfolded simple reflection as a descent, and hence
is the longest unfolded element.
\end{proof}

\subsection{Twist bimodules}

Two-term evaluation bimodules are the standard algebraic model for
spherical twists.  The inverse calculations below follow the method of
Khovanov--Seidel~\cite[Section~2d, Proposition~2.4]{KhovanovSeidel2002};
for fusion-labelled zigzags see
\cite[Definition~5.2, Proposition~5.3]{HengLicata2024}.

All complexes are bounded cochain complexes of finite internally
graded modules or bimodules. Differentials have cochain degree one and
internal degree zero. The internal shift $\langle a\rangle$ raises
degrees by $a$, while $[1]$ lowers cochain placements by one. Tensor
products and Hom complexes use the ordinary cochain signs only; the
internal grading contributes no additional signs. Graded Hom retains
both degrees.

Put $M_\alpha=P_\alpha\otimes_\C Q_\alpha$. Multiplication gives
$\mu_\alpha:M_\alpha\to Z$. The perfect trace pairing between
$P_\alpha$ and $Q_\alpha$ gives the central tensor
\[
 \delta_\alpha(1)=\sum_a p_a\otimes q_a,
 \qquad \tau(q_ap_b)=\delta_{ab},
\]
where the bases are homogeneous. It has internal degree two, and
invariance of the pairing proves its centrality, so it defines a
bimodule map $Z\to M_\alpha\langle-2\rangle$. Define
\begin{equation}
 B_\alpha=[M_\alpha\xrightarrow{\mu_\alpha} Z],
 \qquad
 D_\alpha=[Z\xrightarrow{\delta_\alpha}
                   M_\alpha\langle-2\rangle],
 \label{eq:category-twists}
\end{equation}
with $Z$ in cochain degree zero in both complexes. Thus the projective
terms of $B_\alpha,D_\alpha$ lie in degrees $-1,1$, respectively.

\begin{lemma}
\label{lem:category-twist-inverse}
There are bigraded bimodule chain homotopy equivalences
\[
 B_\alpha\otimes_ZD_\alpha\simeq Z
       \simeq D_\alpha\otimes_ZB_\alpha.
\]
\end{lemma}

\begin{proof}
The product has terms $M_\alpha$ in degree $-1$,
\[
 Z\oplus
 P_\alpha\otimes_\C(e_\alpha Ze_\alpha)
                 \otimes_\C Q_\alpha\langle-2\rangle
\]
in degree zero, and $M_\alpha\langle-2\rangle$ in degree one.
Split the middle factor as
$e_\alpha Ze_\alpha=\C e_\alpha\oplus\C\theta_\alpha$.
The incoming differential projects isomorphically, up to its cochain
sign, onto the token summand. Explicitly, the token coefficient of
$q p_a$ is $\tau(q p_a)$, and the identity
$\sum_a\tau(qp_a)q_a=q$ is the dual-basis identity. Cancel this
degree-$-1$/degree-zero isomorphism. The outgoing differential from
the remaining unit summand is the identity onto
$M_\alpha\langle-2\rangle$, by multiplication with $e_\alpha$;
cancel it as well. The remaining complex is the regular bimodule $Z$
in degree zero. These are Gaussian cancellations of actual bimodule
complexes. In the reverse tensor order the dual-basis identity and
unit multiplication give the same two cancellations, with the
opposite tensor order and the prescribed cochain signs.
\end{proof}

We record the spherical-twist conjugation argument used for braid
relations, following Seidel--Thomas~\cite[Lemma~2.11 and
Propositions~2.12--2.13]{SeidelThomas2001} with the required bimodule
homotopies. For a
bounded perfect left complex $P$, define its evaluation twist to be
the bimodule cone
\begin{equation}
 T_P=\Cone\bigl(P\otimes_\C\Hom_Z(P,Z)\longrightarrow Z\bigr).
 \label{eq:category-evaluation-twist}
\end{equation}
It is unchanged up to homotopy by either shift of $P$. It is also
unchanged under a homotopy equivalence of $P$: a chosen homotopy
inverse and its homotopies identify the two evaluation morphisms up
to homotopy and identify their cones. If $B$ is an invertible
bimodule complex of the kind generated by
\eqref{eq:category-twists}, then
\begin{equation}
 B\otimes_ZT_P\otimes_ZB^{-1}\simeq T_{B\otimes_ZP}.
 \label{eq:category-conjugation}
\end{equation}
Here all terms of $B$ are finite projective on each side: they are
sums of shifts of $Z$ and of $P_\alpha\otimes_\C Q_\beta$.
The left-module dual $B^\vee=\Hom_Z(B,Z)$, with the usual complex
signs, is its tensor adjoint. The evaluation and coevaluation for
finite projectives give the adjunction identities at chain level.
Invertibility therefore identifies $B^\vee$ with $B^{-1}$ in the
bimodule homotopy category. Lemma~\ref{lem:category-frobenius} shows
that these duals remain in the same class of terms. Tensor--Hom
identifies the dual of $B\otimes_ZP$ with
$\Hom_Z(P,Z)\otimes_ZB^\vee$, and under this identification the
composed evaluations agree. Taking their cones proves
\eqref{eq:category-conjugation}.

\begin{lemma}
\label{lem:category-local-braid}
If there is no arrow between distinct unfolded vertices $\alpha,\beta$,
then $B_\alpha B_\beta\simeq B_\beta B_\alpha$. If there is exactly
one arrow in each direction, then
\[
 B_\alpha B_\beta B_\alpha
       \simeq B_\beta B_\alpha B_\beta.
\]
Products here and below denote tensor products over $Z$.
\end{lemma}

\begin{proof}
In the first case $M_\alpha M_\beta=0=M_\beta M_\alpha$, so both
products are the simultaneous evaluation cone with projective term
$M_\alpha\oplus M_\beta$.

In the second case, $B_\alpha P_\beta$ has the nonzero arrow
$P_\alpha\langle1\rangle\to P_\beta$ in degrees $-1,0$.
The complex $D_\beta P_\alpha$ has the nonzero arrow
$P_\alpha\to P_\beta\langle-1\rangle$ in degrees $0,1$.
Nondegeneracy of the edge pairing makes both arrows nonzero, and
their common Hom space is one-dimensional. Therefore
\[
 B_\alpha P_\beta\simeq
       (D_\beta P_\alpha)\langle1\rangle[1],
\]
after rescaling one term if needed. Their evaluation twists agree.
Since $T_{P_\gamma}=B_\gamma$, conjugation gives
\[
 B_\alpha B_\beta D_\alpha
       \simeq D_\beta B_\alpha B_\beta.
\]
Multiplying on the left by $B_\beta$ and on the right by $B_\alpha$
gives the claimed braid relation.
\end{proof}

\begin{proposition}
\label{prop:twist-action}
For each original color $i$, let
\[
 B_i=\prod_{\alpha\in\widetilde{\{i\}}}B_\alpha,
 \qquad
 D_i=\prod_{\alpha\in\widetilde{\{i\}}}D_\alpha.
\]
These products are independent of their order up to bimodule homotopy.
The assignments $\sigma_i\mapsto B_i$ and $\sigma_i^{-1}\mapsto D_i$
define an action of the original Artin group by invertible bimodule
homotopy classes. For $x\in A$ write $B_x$ for a representing complex,
and write $F_i=B_i\otimes_Z-$ on left complexes.

For every $H\in\mathscr C$, right tensoring of labels,
\[
 T_H(P_{s,X})=P_{s,X\otimes H},
\]
is an additive functor on projectives that preserves path lengths and
extends to their bounded homotopy category. It commutes naturally up
to homotopy with the positive and inverse color twists, and hence
with each $B_x\otimes_Z-$. All the action identities above are
bimodule chain homotopy equivalences, so they remain valid upon
tensoring with finite complexes of modules that are not projective.
\end{proposition}

\begin{proof}
Distinct vertices in a color have no mutual Hom spaces. Their twists
commute by Lemma~\ref{lem:category-local-braid}, and their products
are inverse by Lemma~\ref{lem:category-twist-inverse}. Moreover
\begin{equation}
 B_i\cong
 \Cone\left(\bigoplus_{\alpha\in\widetilde{\{i\}}}
               P_\alpha\otimes_\C Q_\alpha\longrightarrow Z\right):
 \label{eq:category-color-cone}
\end{equation}
every mixed tensor term in the expansion vanishes by orthogonality.

If $m_{ij}<\infty$, the subsystem on $\{i,j\}$ is spherical.
The two alternating words of length $m_{ij}$ are reduced expressions
for its longest element. By Lemma~\ref{lem:finite-unfolding}, both
lift to reduced expressions for the same longest element of its
finite simply-laced unfolding. Matsumoto's theorem relates these
expressions by commutations and three-term braid moves, all of which
hold for the bimodules by Lemma~\ref{lem:category-local-braid}.
This proves the original Artin relation. When $m_{ij}=\infty$ there
is no relation to verify. Thus the assignment defines the asserted
homotopy action for every Coxeter matrix under consideration.

On a morphism in \eqref{eq:category-projective-hom}, $T_H$ is
right tensoring by $\id_H$, with the associativity identifications.
All cap multiplications are to the left of the label, so this
preserves composition and path length. In the evaluation term of
\eqref{eq:category-color-cone} applied to $P_{t,Y}$, semisimplicity
gives the canonical identification
\begin{align*}
 &\bigoplus_{a,L}
    P_{i,L}\langle a\rangle\otimes_\C
      \Hom_{\mathscr C}(L,D^a_{it}\otimes Y)\\
 &\hspace{35mm}\cong
      \bigoplus_a P_{i,D^a_{it}\otimes Y}\langle a\rangle.
\end{align*}
Its evaluation is the tautological label map. Tensoring this
description on the right by $H$ is the same description with $Y$
replaced by $Y\otimes H$, naturally in $Y$ and its projective
morphisms. Hence the evaluation cones commute naturally with $T_H$
on projectives and on their complexes. Whiskering this natural
commutation by the inverse equivalence proves the same assertion
for $D_i$, without requiring $T_H$ itself to be invertible.

Finally, tensoring an explicit chain homotopy and its equations with
any module complex preserves those equations with the ordinary
cochain tensor signs. The last assertion therefore follows from
the homotopy-level inverse and braid calculations, without any
flatness requirement on these additional module complexes.
\end{proof}

We will extract the height data from the images $B_xP_{o,L}$ of the
frame projectives, for $x\in A$ and all simple labels $L$.  Retaining
the whole label family will allow us to combine its generator counts,
using the positive label dimensions, into vectors indexed by the
original colors and the frame.  We first need to extract these counts
from homotopy classes.  The next section uses minimal complexes and
groups their generators into layers according to $c+u$, where $c$ is
the cochain placement and $u$ is the internal generator degree.
